\chapter{Introduction}
\label{intro}

Algebraic geometry is the study of solutions to polynomial equations. These solutions can be studied both from an algebraic and geometric perspective, and the interplay between these two viewpoints is what makes algebraic geometry such a rich and interesting field of mathematics.

In 1849, Arthur Cayley communicates to George Salmon that a general cubic surface contains a finite number of lines, which Salmon proves must be equal to 27. Later this same year, Salmon proves that in fact any non-singular surface contains exactly 27 lines. \cite{Do04} claimed it can be considered as the beginning of modern algebraic geometry.


\subsection{Outline}
We prove that a general cubic surface contains exactly 27 lines, building up to this result by introducing necessary background theory. This includes the theory of Chern classes, the universal Fano variety, the Grassmannian. We then introduce the notion of $m$-crosses, $m$-gons and $m$-chains of lines on surfaces, and study when surfaces contain such configurations of lines. We include computational results for these configurations.